% Source printed pp. 513-514 / PDF pp. 521-522.
\subsection{Bridgeland 稳定性}\label{src:10.2.4-subsection}
Bridgeland 引入了光滑射影簇 $X$ 上凝聚层的有界导出范畴 $\mathcal D^b(X)$ 上的\newterm{稳定性条件}{stability condition}的概念 \cite{Bri07}.\footnote{Bridgeland 更一般地考虑任意三角范畴上的稳定性条件.} 这个概念近年来引发了大量研究, 促成了许多新的模空间的构造, 以及经典代数几何中若干未解决问题的解答 \cite{BM23}.

$D^b(X)$ 上的稳定性条件 $\sigma=(Z,\mathcal A)$ 由群同态 $Z:K_{\mathrm{num}}(X)\to\mathbb C$ 以及有界 $t$-结构的\newterm{心}{heart} $\mathcal A\subseteq D^b(X)$ 组成, 前者称为\newterm{中心荷}{central charge}. $\mathcal A$ 中的复形 $E$ 称为 $\sigma$-半稳定, 若对 $\mathcal A$ 中每个非零子对象 $F\subseteq E$ 都有 $\arg Z(F)\leq\arg Z(E)$; 任意复形 $E\in D^b(X)$ 称为 $\sigma$-半稳定, 若对某个 $k\in\mathbb Z$, $E[k]\in\mathcal A$ 是 $\sigma$-半稳定的. 要成为稳定性条件, $\sigma=(Z,\mathcal A)$ 必须满足若干性质:
\begin{enumerate}[(1)]
\item 对非零 $E\in\mathcal A$, 荷 $Z(E)$ 位于严格上半平面或严格负实轴;
\item 每个对象 $E$ 都有 Harder--Narasimhan 滤过, 即存在有限复形态射序列 $0=E_0\to E_1\to\cdots\to E_n=E$, 使 $E_i\to E_{i+1}$ 的锥 $\sigma$-半稳定且相位递减;
\item 满足\newterm{支撑性质}{support property}. 粗略地说, 它要求存在 $\Lambda_{\mathbb R}$ 上的二次型 $Q$, 使 $\ker(Z)$ 关于 $Q$ 负定, 且对每个 $\sigma$-半稳定对象 $E$ 都有 $Q(E)\geq0$.
\end{enumerate}

这一理论的一个突出特点是, 稳定性条件空间 $\mathrm{Stab}(X)$ 自然成为复流形, 并且将 $\sigma=(Z,\mathcal A)$ 映为 $Z$ 的遗忘映射 $\mathrm{Stab}(X)\to\mathrm{Hom}(K_{\mathrm{num}}(X),\mathbb C)$ 是局部同构 \cite{Bri07}. 对亏格 $g\geq1$ 的光滑连通射影曲线 $C$, $D^b(C)$ 上有自然稳定性条件 $\sigma=(Z,\mathrm{Coh}(C))$, 其中 $Z:K_{\mathrm{num}}(C)\to\mathbb C$ 定义为 $Z(E)=-\deg(E)+i\operatorname{rk}(E)$, $\mathrm{Coh}(C)$ 是标准的心. 在行列式为正的矩阵群的万有覆盖 $\widetilde{\mathrm{GL}}_2^+(\mathbb R)$ 的作用下, 这给出了 $D^b(C)$ 上唯一的稳定性条件.

固定 $D^b(X)$ 上的稳定性条件 $\sigma=(Z,\mathcal A)$, 有模叠 $\mathcal M_{\mathcal A}$ 参数化 $\mathcal A$ 中对象的平坦族. 这个叠可抽象地定义为 Abel 范畴 $\mathcal A$ 中对象的模叠, 其中 $\Bbbk$-代数 $R$ 上的对象族是 $\mathcal A$ 中有限表现且平坦的 $R$-模对象 (即复形 $E\in\mathcal A$ 连同 $\Bbbk$-代数映射 $R\to\mathrm{End}_{\mathcal A}(E)$) \cite{AZ01,AP06,Pol07}. 另一种观点是将 $\mathcal M_{\mathcal A}$ 视为 Lieblich 的代数叠的开子叠; 后者参数化负次数自 Ext 消失的复形 \cite{Lie06a}, \cite[Tag 0DLN]{SP}. 叠 $\mathcal M_{\mathcal A}$ 局部有限型且具有仿射对角态射. 固定类 $v\in K_{\mathrm{num}}(X)$, 考虑参数化类为 $v$ 的对象的子叠 $\mathcal M_{\mathcal A}(v)$.

一般而言, 尚不清楚 $\sigma$-半稳定对象的子叠 $\mathcal M_{\mathcal A}^{\sigma\text{-}\mathrm{ss}}(v)$ 是否开放且有界. 不过 Toda 证明了, 若相对于 $\sigma$ 的半稳定轨迹开放且有界, 则对与 $\sigma$ 位于同一连通分支中的每个稳定性条件 $\sigma'\in\mathrm{Stab}(X)$, 相应的半稳定轨迹也开放且有界 \cite{Tod08,PT19}. 在某些情形, 例如光滑射影曲面上, 存在与 Gieseker 稳定性等同的稳定性条件, 其半稳定轨迹已知开放且有界. 假设相对于 $\sigma$ 的半稳定性开放且有界, 则存在定理 (\bookref{7.10.1}) 给出固有良模空间 $\mathcal M_{\mathcal A}^{\sigma\text{-}\mathrm{ss}}(v)\to M_{\mathcal A}^{\sigma\text{-}\mathrm{ps}}(v)$ \cite[\S7]{AHLH23}, \cite{BLM+21}. 射影性在一般情形仍未知, 但已在以下情形建立: $X$ 为 K3 曲面且 $\sigma$ 一般 \cite{BM14}; $X$ 为光滑曲面且 $\sigma$ 位于``竖直壁''上 \cite{Taj23}; 以及若干其他情形, 例如模问题有 GIT 描述时 \cite{ABCH13,AM17}.
